% Generated draft: requires mathematical and source review.
\documentclass[11pt]{article}
\usepackage[margin=25mm]{geometry}
\usepackage{fontspec}
\setmainfont{Latin Modern Roman}
\usepackage{amsmath,amssymb,mathtools,unicode-math}
\setmathfont{Latin Modern Math}
\usepackage{enumitem,needspace,xurl,hyperref}
\setlength{\parindent}{0pt}
\setlength{\parskip}{5pt}
\setlength{\emergencystretch}{3em}
\newcommand{\N}{\mathbb{N}}
\newcommand{\Z}{\mathbb{Z}}
\newcommand{\Q}{\mathbb{Q}}
\newcommand{\R}{\mathbb{R}}
\newcommand{\C}{\mathbb{C}}
\newcommand{\F}{\mathbb{F}}
\newcommand{\A}{\mathbb{A}}
\newcommand{\PP}{\mathbb{P}}
\newcommand{\E}{\mathbb{E}}
\newcommand{\eps}{\varepsilon}
\newcommand{\dd}{\,\mathrm{d}}
\begin{document}
\section*{UnsolvedMath ID: 30005399}
\textbf{Source problem code:} catalogue entry
\subsection{Definitions and mathematical statement}
Let \(p\) denote the prime occurring in the supplied record, and write \(X=X_{E,F}\) for its Fargues--Fontaine curve. The record does not specify the hypotheses on the parameters \(E\) and \(F\). Its extracted theorem distinguishes an algebraic curve \(X^{\mathrm{alg}}\) and an adic space \(X^{\mathrm{ad}}\), but does not supply their precise constructions or comparison conventions. We write \(\mathcal G\) for a coefficient sheaf, to distinguish it from the parameter \(F\).

The requested problem is to describe a dualizing complex for \(X\), and to determine the associated cohomological behavior with \(p\)-torsion coefficients. This is an unresolved formulation rather than a uniquely specified conjecture: the record does not identify the duality theory, coefficient ring, site, structural morphism, or normalization of the proposed dualizing complex. In particular, it does not specify whether coherent duality or a duality theory for torsion sheaves is intended. It also does not specify whether the general phrase \(p\)-torsion means sheaves annihilated by \(p\), sheaves annihilated by a power of \(p\), or sheaves whose local sections are each annihilated by some power of \(p\).

The following supplied background qualifies the problem and must be retained.
\begin{enumerate}
\item The extracted report states that, for an \(\acute{e}\)tale abelian torsion sheaf \(\mathcal G\) on \(X^{\mathrm{alg}}\) whose torsion is prime to \(p\), its Conjectures \((1)\) and \((2)\) hold. Those conjectures are not included in the record, so their mathematical content cannot be restated. The extract additionally reports
\[
\operatorname{cd}_{\ell}\!\left(E(X^{\mathrm{alg}})\right)\leq 1,
\qquad \ell\neq p\text{ a prime}.
\]
Here \(\operatorname{cd}_{\ell}\) denotes \(\ell\)-cohomological dimension; the field notation \(E(X^{\mathrm{alg}})\) is retained from the extract, whose precise conventions are missing. These are reported background assertions, not conclusions established by this draft.
\item For a constructible \(\mathbf F_p\)-module sheaf \(\mathcal G\), the extract reports the conditional assertion
\[
\text{(Stein)}\quad\Longrightarrow\quad
H^i\!\left(X^{\mathrm{ad}},\mathcal G^{\mathrm{ad}}\right)=0
\quad\text{for every integer }i\geq 3.
\]
A \(\mathbf F_p\)-module sheaf is, in particular, annihilated by \(p\). The notation \(\mathcal G^{\mathrm{ad}}\) represents the associated sheaf appearing in the extract; its construction and the precise cohomological site are not specified there.
\item The hypothesis labelled \(\text{(Stein)}\) concerns the tilt of a perfectoid space obtained by cyclotomic base change of an adification of an open subset. The extract describes \(E_\infty/E\) as the cyclotomic \(\mathbf Z_p\)-extension and asks that the resulting tilt be a perfectoid Stein space. However, its description of the open subset mixes algebraic and adic notation, its completed-field notation is corrupted, and the purported definition ends with an incomplete union of spaces \(V_n\). The conditions on \(V_n\), the exhaustion, and the required transition maps are absent. Consequently, \(\text{(Stein)}\) is only a label for a missing full hypothesis here, and the displayed implication is a record of the supplied conditional assertion rather than a self-contained theorem.
\end{enumerate}

Thus the source supports the research questions above and the reported conditional vanishing in degrees at least \(3\). It does not specify an unconditional vanishing conjecture for all \(p\)-torsion coefficients, a complete duality statement, or a precise relation between those two targets. These missing qualifications require human mathematical review.
\subsection{Short English statement}
Describe a dualizing complex for the Fargues--Fontaine curve and clarify its cohomology with \(p\)-torsion coefficients. The supplied extract reports vanishing in degrees \(i\geq 3\) for constructible \(\mathbf F_p\)-coefficients under a Stein hypothesis, but that hypothesis is truncated. The curve parameters, duality framework, and scope of \(p\)-torsion coefficients remain unspecified.
\subsection{Sources}
\begin{itemize}
\item Supplied local registry record: \texttt{data/unsolvedmath/problems.json}, permanent UnsolvedMath ID \(30005399\), source problem code \texttt{catalogue entry}, title \emph{Dualizing Complexes and p-Torsion on the Fargues--Fontaine Curve}. No individual proposer is identified; the supplied proposed year is \(2023\).
\item Supplied report citation: \emph{Arithmetic of Shimura Varieties} (2023), Oberwolfach Reports, pp.\ 261--326, DOI \texttt{10.4171/owr/2023/5}. \url{https://doi.org/10.4171/owr/2023/5}. The record attributes the problem to workshop open-problem discussions documented in this report; it gives no precise page or speaker for this item.
\item Supplied literature source, labelled ``OWR-2023-5, open questions'': \url{https://oa.tib.eu/renate/bitstreams/ae468b44-47d2-4568-9d88-ffcb4023f8d4/download}.
\item Supplied literature source, labelled ``Bartling, partial cohomology results'': \url{https://arxiv.org/abs/2206.14253}. The extract refers to a talk's main result as reference [1], but does not provide the corresponding bibliography entry; its identification with this link is not established by the supplied text.
\item Literature assessment checked at \texttt{2026-08-22}: the supplied curation labels the problem open and describes the dualizing-complex and \(p\)-torsion questions as speculative, stating that existing Banach--Colmez and \(\acute{e}\)tale-cohomology results do not provide a general answer. This is supplied dated curation, not an independently verified literature assessment. The record itself requests independent verification of time-sensitive status and cited claims.
\item The registry labels the statement \texttt{corrected\_verified} and says it was reconstructed and checked against the official Oberwolfach report. That is a supplied provenance claim; no source visit or independent verification is asserted in this draft.
\end{itemize}
\end{document}
